\slajdid{10-s089}
\begin{frame}[label=10-s089]
\gusttekst\setlength{\abovedisplayskip}{3pt}\setlength{\belowdisplayskip}{3pt}\setlength{\abovedisplayshortskip}{2pt}\setlength{\belowdisplayshortskip}{2pt}\begin{columns}[T,onlytextwidth]\column{.41\textwidth}\centering\begin{tikzpicture}[x=.45cm,y=.55cm,font=\sitneoznake,>=Latex]\ctikzset{bipoles/length=.65cm}
\draw(-1.9,3)--(1.9,3);\node[above]at(0,3){$V_{CC}$};
\draw(-2,-.5)--(-2,.5);\draw(-2,.32)--(-1.3,.7)--(-1.3,1.3)to[R,l=$R_{C1}$](-1.3,3);\draw[->](-2,-.32)--(-1.3,-.7);\draw(-1.3,-.7)--(-1.3,-1.1)--(1.3,-1.1);\draw(-2,0)--(-3.1,0)node[ocirc]{}node[left]{$V_I$};\node[right]at(-1.7,0){T1};\draw(-1.3,1.1)--(-2.5,1.1)node[ocirc]{}node[left]{$V_{O1}$};
\draw(2,-.5)--(2,.5);\draw(2,.32)--(1.3,.7)--(1.3,1.3)to[R,l_=$R_{C2}$](1.3,3);\draw[->](2,-.32)--(1.3,-.7);\draw(1.3,-.7)--(1.3,-1.1);\draw(2,0)--(3.1,0)node[ocirc]{}node[right]{$V_R$};\node[left]at(1.7,0){T2};\draw(1.3,1.1)--(2.5,1.1)node[ocirc]{}node[right]{$V_{O2}$};
\node[draw,circle,minimum size=7mm,inner sep=0pt](ie)at(0,-2){};\draw[->,shorten >=2pt,shorten <=2pt](ie.north)--(ie.south);\draw(0,-1.1)--(ie.north);\draw(ie.south)--(0,-3);\draw(-.65,-3)--(.65,-3);\node[below]at(0,-3){$V_{EE}$};\node[right]at(.6,-2){$I_E$};\end{tikzpicture}
\hfill\raisebox{1.7cm}{\tikz\draw[->,DEplava](0,0)--(.8,0);}\par\medskip\raggedright
Ovde treba uočiti da će bilo kakva struja kroz izlaz poremetiti ove uslove, odnosno da se lako može desiti da tranzistori odu u zasićenje zbog izlaznih struja. Zbog toga se u realnom ECL kolu dodaju izlazni baferi koji će obezbediti povećane strujne kapacitete ali neće poremetiti naponske nivoe (pojačavač sa zajedničkim kolektorom)
\column{.57\textwidth}\centering\begin{tikzpicture}[x=.5cm,y=.48cm,font=\sitneoznake,>=Latex]\ctikzset{bipoles/length=.65cm}
\draw(-4.8,3)--(4.8,3);\node[above]at(0,3){$V_{CC}$};
\draw(-2,-.5)--(-2,.5);\draw(-2,.32)--(-1.3,.7)--(-1.3,1.3)to[R,l=$R_{C1}$](-1.3,3);\draw[->](-2,-.32)--(-1.3,-.7);\draw(-1.3,-.7)--(-1.3,-1.1)--(1.3,-1.1);\draw(-2,0)--(-3.1,0)node[ocirc]{}node[left]{$V_I$};\node[right]at(-1.7,0){T1};
\draw(2,-.5)--(2,.5);\draw(2,.32)--(1.3,.7)--(1.3,1.3)to[R,l_=$R_{C2}$](1.3,3);\draw[->](2,-.32)--(1.3,-.7);\draw(1.3,-.7)--(1.3,-1.1);\draw(2,0)--(3.1,0)node[ocirc]{}node[right]{$V_R$};\node[left]at(1.7,0){T2};
\node[draw,circle,minimum size=7mm,inner sep=0pt](ie)at(0,-2){};\draw[->,shorten >=2pt,shorten <=2pt](ie.north)--(ie.south);\draw(0,-1.1)--(ie.north);\draw(ie.south)--(0,-3);\draw(-.65,-3)--(.65,-3);\node[below]at(0,-3){$V_{EE}$};\node[right]at(.6,-2){$I_E$};\draw(-3.8,.6)--(-3.8,1.6);\draw(-1.3,1.1)--(-3.8,1.1);\draw(-3.8,1.42)--(-4.5,1.8)--(-4.5,3);\draw[->](-3.8,.78)--(-4.5,.4);\draw(-4.5,.4)--(-4.5,.1)--(-5.7,.1)node[ocirc]{}node[left]{$V_{O1}$};\node[left]at(-4,1.1){T3};
\draw(3.8,.6)--(3.8,1.6);\draw(1.3,1.1)--(3.8,1.1);\draw(3.8,1.42)--(4.5,1.8)--(4.5,3);\draw[->](3.8,.78)--(4.5,.4);\draw(4.5,.4)--(4.5,.1)--(5.7,.1)node[ocirc]{}node[right]{$V_{O2}$};\node[right]at(4,1.1){T4};\node at(-6,-2.8){$\cdots$};\end{tikzpicture}
\par\medskip\raggedright
Tranzistori T3 i T4 će zbog svoje konfiguracije, podrazumevajući da postoji neka struja na izlazu, uvek raditi u aktivnom režimu. Menjaju se naponi logičke nule i jedinice
\[V_{OL}=V_{CC}-R_C I_E-V_{BE}\]
\[V_{OH}=V_{CC}-V_{BE}\]
Uslovi za izbor elemenata se ne menjaju.
\end{columns}
\end{frame}
